
\section{Low-rank matrix denoising}
\label{sec:matrix-denoising-L2}

To catch a glimpse of the effectiveness of the approach we have introduced,
let us start by trying it out on a warm-up example: low-rank matrix denoising.

\subsection{Problem formulation and algorithm}\label{sec:matrix-denoising-setup}  Consider an unknown rank-$r$ symmetric matrix $\bm{M}^{\star} \in \mathbb{R}^{n\times n}$ with eigendecomposition $\bm{M}^{\star}=\bm{U}^{\star}\bm{\Lambda}^{\star}\bm{U}^{\star\top}$, where the columns of $\bm{U}^{\star}\in \mathbb{R}^{n\times r}$ are orthonormal, and $\bm{\Lambda}^{\star}\in \mathbb{R}^{r\times r}$ is a diagonal matrix containing the nonzero eigenvalues $\{\lambda_{i}^{\star}\}$ of $\bm{M}^{\star}$. Assume that $|\lambda_1^{\star}|\geq |\lambda_2^{\star}|\geq \cdots \geq |\lambda_r^{\star}|>0$.  Suppose that  we observe a noisy copy
\[
\bm{M}=\bm{M}^{\star}+\bm{E},
\]
 where $\bm{E}=[E_{i,j}]_{1\leq i,j\leq n}$ is a symmetric noise matrix. It is assumed that the entries $\{{E}_{i,j}\}_{i\geq j}$ are independently generated obeying
%
\begin{align}
	E_{i,j} \overset{\mathrm{i.i.d.}}{\sim} \mathcal{N}(0,\sigma^2) , \qquad i\geq j.
\end{align}
%
The aim is to estimate the eigenspace $\bm{U}^{\star}$ from the data matrix $\bm{M}$. Despite its simplicity, this problem  has been extensively studied in the literature \citep{koltchinskii2016perturbation,bao2018singular,ding2020high,xia2019normal,li2021minimax}. It also  bears close relevance to the famous angular/phase synchronization problem \citep{singer2011angular,bandeira2017tightness}.


In order to estimate the low-rank factors specified by $\bm{U}^{\star}$, a natural scheme is to resort to the rank-$r$ leading eigenspace of the data matrix $\bm{M}$.  More precisely, denote by $\lambda_1,\cdots,\lambda_n$ the eigenvalues of $\bm{M}$ sorted by their magnitudes, i.e.,
%
\begin{align}
	|\lambda_1| \geq |\lambda_2| \geq \cdots \geq |\lambda_n|,
\end{align}
%
and let $\bm{u}_1$, $\cdots$, $\bm{u}_n$ represent the associated eigenvectors.
This spectral method returns  $\bm{U} = [\bm{u}_1,\cdots,\bm{u}_r]\in \mathbb{R}^{n\times r}$ as an estimate of $\bm{U}^{\star}$.


\subsection{Performance guarantees}
\label{sec:performance-denoising-L2}

\paragraph{Statistical accuracy of the spectral estimate.}

We now examine the accuracy of the above spectral estimate. Towards this,
a key step lies in bounding the spectral norm of the noise matrix
$\bm{E}$. We claim for the moment that (which will be established in Section~\ref{sec:proof-eqn-noise-matrix-E-bound-denoising})
%
\begin{align}
	\|\bm{E}\| \leq 5\sigma\sqrt{n}
	\label{eq:noise-matrix-E-bound-denoising}
\end{align}
%
with probability at least $1-O(n^{-8})$.
Armed with this claim and the fact $\lambda_{r+1}^{\star}=0$,
we are in a situation where it is quite easy to see how the Davis-Kahan theorem applies.
According to Corollary \ref{cor:davis-kahan-conclusion-corollary}, with probability greater than $1-O(n^{-8})$ one has
%
\begin{align}
	\mathsf{dist}\big(\bm{U},\bm{U}^{\star}\big)\leq\frac{2\|\bm{E}\|}{|\lambda_{r}^{\star}|}\leq\frac{10\sigma\sqrt{n}}{|\lambda_{r}^{\star}|} ,
	\label{eq:dist-U-Ustar-denoising-L2}
\end{align}
%
provided that the noise variance is sufficiently small obeying $ \sigma \sqrt{n} \leq \frac{1-1/\sqrt{2}}{5} |\lambda_r^{\star}|$ so that $\|\bm{E}\|\leq (1-1/\sqrt{2})|\lambda_r^{\star}|$.




\paragraph{Tightness and optimality.}
%
The tightness of the statistical guarantee~\eqref{eq:dist-U-Ustar-denoising-L2} can be assessed when compared with the minimax lower bound. For instance, it is well-known in the literature (e.g., \citet[Theorem 3]{cheng2020tackling}) that: even for the case with $r=1$, one cannot hope to achieve $\mathsf{dist}\big(\widehat{\bm{U}},\bm{U}^{\star}\big)=o(\sigma \sqrt{n}/|\lambda_{r}^{\star}| )$---in a minimax sense---regardless of the estimator $\widehat{\bm{U}}$ in use. 
Consequently, the spectral method turns out to be orderwise statistically optimal for low-rank matrix denoising.




%(in addition to the proof of Corollary \ref{cor:davis-kahan-conclusion-corollary})

\paragraph{Additional useful results: eigenvalue and matrix estimation.}
	Before concluding, we record several immediate consequences of the above analysis that will be useful later on. Specifically, assuming that $ \sigma \sqrt{n} \leq \frac{1-1/\sqrt{2}}{5} |\lambda_r^{\star}|$, we see from Weyl's inequality (cf.~Lemma \ref{lemma:weyl}) that
	%
	\begin{align}
		|\lambda_{i}| & \leq\|\bm{E}\|\leq5\sigma\sqrt{n},
		\qquad
		\text{for all }i\geq r+1
		\label{eq:lambda-i-bound-E-denoising}
	\end{align}
	%
	% and \cm{This is wrong.}
	% %
	% \begin{align}
	% 	|\lambda_{i} - \lambda_i^{\star}| & \leq\|\bm{E}\|\leq5\sigma\sqrt{n}, \qquad
	% 	\text{for all }i\leq r
	% 	\label{eq:lambda-1-bound-E-denoising}
	% \end{align}
	%
with probability $1-O(n^{-8})$.

We further remark on the Euclidean statistical accuracy when estimating the unknown matrix $\bm{M}^{\star}$ using $\widehat{\bm{M}} \coloneqq \bm{U}\bm{\Lambda}\bm{U}^{\top}$, where $\bm{\Lambda} \coloneqq \mathsf{diag}\big([\lambda_{1},\cdots,\lambda_{r}]\big)$.
It is seen from the triangle inequality that
%
\begin{align}
\big\|\bm{U}\bm{\Lambda}\bm{U}^{\top}-\bm{M}^{\star}\big\| & \leq\big\|\bm{M}-\bm{M}^{\star}\big\|+\big\|\bm{U}\bm{\Lambda}\bm{U}^{\top}-\bm{M}\big\|\nonumber \\
 & =\big\|\bm{E}\big\|+\big|\lambda_{r+1}\big|\leq2\big\|\bm{E}\big\|,
	\label{eq:ULambdaU-Mstar-norm-denoising}
\end{align}
%
where the last inequality relies on \eqref{eq:lambda-i-bound-E-denoising}.
Since the rank of $\bm{U}\bm{\Lambda}\bm{U}^{\top}-\bm{M}^{\star}$ is at most $2r$,
with probability at least $1-O(n^{-8})$ one has
%
\begin{align}
	\big\|\bm{U}\bm{\Lambda}\bm{U}^{\top}-\bm{M}^{\star}\big\|_{\mathrm{F}}
	& \leq\sqrt{2r} \, \big\|\bm{U}\bm{\Lambda}\bm{U}^{\top}-\bm{M}^{\star}\big\|  \leq 2\sqrt{2r}\, \|\bm{E}\| \nonumber \\
	& \leq 10 \sigma\sqrt{2nr}.
	\label{eq:top-r-approx-Fro-bound}
\end{align}


\subsection{Proof of the inequality \eqref{eq:noise-matrix-E-bound-denoising} on $\|\bm{E}\|$}
\label{sec:proof-eqn-noise-matrix-E-bound-denoising}

We plan to employ Theorem~\ref{thm:tighter-spectral-normal-ramon}.
Given that Gaussian entries are unbounded, we introduce a truncated copy $\widetilde{\bm{E}}=[\widetilde{E}_{i,j}]_{1\leq i,j\leq n}$ defined as follows
%
\begin{equation}
	\widetilde{E}_{i,j}\coloneqq E_{i,j} \mathbbm{1}\big\{|E_{i,j}|\leq 5\sigma\sqrt{\log n}\big\}, \quad 1\leq i,j\leq n.
	\label{eq:defn-Eij-truncated-denoising}
\end{equation}
%
Two properties are in place.
\begin{itemize}
\item It is readily seen from the property of Gaussian distributions
that
\[
\mathbb{P}\big\{ E_{i,j}=\widetilde{E}_{i,j}\big\}\geq1-n^{-12},\qquad1\leq i,j\leq n,
\]
which combined with the union bound leads to
\begin{equation}
\mathbb{P}\big\{\bm{E}=\widetilde{\bm{E}}\big\}\geq1-n^{-10}.\label{eq:prob-E-Etilde-equivalence-denoising}
\end{equation}
\item Given that $B\coloneqq\max_{i,j}|\widetilde{E}_{i,j}|\leq 5\sigma\sqrt{\log n}$,
we can invoke Theorem~\ref{thm:tighter-spectral-normal-ramon} (or more directly,
(\ref{eq:X-spectral-norm-iid-special-ramon})) to demonstrate that
\[
\|\widetilde{\bm{E}}\|\leq4\sigma\sqrt{n}+O(B\log n)\leq5\sigma\sqrt{n}
\]
for sufficiently large $n$, 
with probability exceeding $1-O(n^{-8})$. Here, we implicitly use the fact that $\mathbb{E}[\widetilde{{E}}_{i,j}^2] \leq \mathbb{E}[E_{i,j}^2 ] = \sigma^2$.
\end{itemize}
Combining the above two observations implies that
%
\[
\|\bm{E}\|=\|\widetilde{\bm{E}}\|\leq5\sigma\sqrt{n}
\]
%
 with probability exceeding $1-O(n^{-8})$, as claimed.
 %\end{proof}



%\begin{remark}
%	Before concluding, we record several immediate consequences of the above analysis (in addition to the proof of Corollary \ref{cor:davis-kahan-conclusion-corollary}) that will prove useful later on. Specifically, assuming that $ \sigma \sqrt{n} \leq \frac{1-1/\sqrt{2}}{5} |\lambda_r^{\star}|$ , we have --- with probability $1-O(n^{-8})$ --- that
%	%
%	\begin{align}
%		|\lambda_{i}| & \leq\|\bm{E}\|\leq5\sigma\sqrt{\log n},
%		\qquad
%		\text{for all }i\geq r+1.
%		\label{eq:lambda-i-bound-E-denoising}
%	\end{align}
%	%
%	In the special case where $r=1$, with probability $1-O(n^{-8})$ one has
%	%
%	\begin{align}
%		|\lambda_{1} - \lambda_1^{\star}| & \leq\|\bm{E}\|\leq5\sigma\sqrt{\log n}.
%		\label{eq:lambda-1-bound-E-denoising}
%	\end{align}
%	%
%\end{remark}


